\documentclass[11pt]{article}
\usepackage[margin=1in]{geometry}
\usepackage{amsmath,amssymb,amsthm,mathtools,booktabs,microtype}
\usepackage[dvipsnames]{xcolor}
\usepackage{array}
\usepackage[colorlinks=true,linkcolor=blue!50!black,citecolor=blue!50!black]{hyperref}

% --- project macros (casper_k_macros.tex) ---
\providecommand{\casper}{\textnormal{\textsc{Casper}}}
\providecommand{\R}{\mathbb{R}}
\DeclareMathOperator{\Topk}{top}\renewcommand{\Topk}{\operatorname{top}_k}
\newcommand{\topkS}{\operatorname{top}_k^S}
\DeclareMathOperator{\supp}{supp}
% --- local ---
\newcommand{\yh}{\hat y}
\newcommand{\zt}{\tilde z}
\newcommand{\zm}{z^{\mathrm{mask}}}
\newcommand{\Sb}{\bar S}
\newcommand{\1}{\mathbf 1}
\newcommand{\Ck}{\mathcal C_k^{\epsilon}}
\newcommand{\Rc}{R_{\casper}}
\newcommand{\Rm}{R_{\mathrm{mask}}}
\newcommand{\Ent}{\mathrm{H}}
\DeclareMathOperator*{\argmin}{arg\,min}

\newtheorem{proposition}{Proposition}
\newtheorem{lemma}[proposition]{Lemma}
\newtheorem{corollary}[proposition]{Corollary}
\theoremstyle{definition}
\newtheorem{definition}[proposition]{Definition}
\theoremstyle{remark}
\newtheorem{remark}[proposition]{Remark}

\title{Why $\casper_k$'s restoration is a better-aimed training signal\\ than a forbidden-mass penalty}
\author{CASPER-$k$ project notes}
\date{\today}

\begin{document}
\maketitle

\begin{abstract}
Masking a policy's output distribution guarantees compliant actions but passes no gradient to the
forbidden actions, so the network never learns the constraint. Two training signals restore one:
a \emph{penalty} on the forbidden probability mass, and $\casper_k$'s \emph{restoration} term, the
squared distance to the $\casper_k$ projection. We show that (i) restoration vanishes, together with
its gradient, exactly on the set of $\epsilon$-separated $k$-compliant distributions, while the
penalty vanishes nowhere; consequently (ii) every compliant task-optimum remains stationary under
restoration for any weight, and under the penalty for no positive weight; (iii) the penalty step
provably lowers the entropy of the executed (masked) policy at every state where a forbidden action
has mass; and (iv) restoration's gradient is supported on the actions that compete for the top-$k$
slots and preserves mass. The same properties predict when the advantage should disappear: for
restoration towards the masked distribution, and for $\casper_k$'s face projection when $|S|<k$.
Both predictions match the NAVIX experiments.
\end{abstract}

\section{Setting}
\label{sec:setting}

Consider one categorical decision with $h$ values, logits $\ell\in\R^h$ and output
$\yh=\operatorname{softmax}(\ell)\in\Delta^\circ$, the open simplex. Let $S\subseteq[h]$ be the
allowed set, $\Sb=[h]\setminus S\neq\emptyset$, and
\[
F(\yh)=\sum_{j\in\Sb}\yh_j\in(0,1),\qquad P_S=1-F .
\]
In RL, $\yh=\yh_\theta(\cdot\mid s)$ and $S=S(s)$ depend on the state; in states with $S=[h]$ all
regularisers below are identically zero.

\paragraph{Masking.} The masked distribution and its top-$k$ sampler are
\[
\zm_a=\frac{\yh_a\,\1[a\in S]}{P_S},\qquad
T^\star=\topkS(\yh),\qquad
q^{\mathrm{mask}}_a=\frac{\zm_a\,\1[a\in T^\star]}{\sum_{b\in T^\star}\zm_b}.
\]

\begin{lemma}[Masking discards the forbidden logits]\label{lem:mask}
For every $a\in S$ and $j\in\Sb$, $\partial\log\zm_a/\partial\ell_j=0$ and
$\partial\log q^{\mathrm{mask}}_a/\partial\ell_j=0$. Hence any loss that depends on $\yh$ only
through $\zm$ or $q^{\mathrm{mask}}$ (in particular the policy-gradient loss of a masked policy) has
zero gradient with respect to every forbidden logit.
\end{lemma}
\begin{proof}
$\log\zm_a=\ell_a-\log\sum_{b\in S}e^{\ell_b}$ contains no $\ell_j$ with $j\in\Sb$; the same holds
for $\log q^{\mathrm{mask}}_a=\ell_a-\log\sum_{b\in T^\star}e^{\ell_b}$, since $T^\star\subseteq S$
is determined by the ranking of the allowed logits only.
\end{proof}

\paragraph{The two training signals.} With task loss $\mathcal L_{\mathrm{task}}$, we compare
\begin{align}
\mathcal J_{\mathrm{pen}}&=\mathcal L_{\mathrm{task}}+\lambda\,\mathbb E_s\bigl[F(\yh)\bigr],
\label{eq:pen}\\
\mathcal J_{\casper}&=\mathcal L_{\mathrm{task}}+\lambda_r\,\mathbb E_s\bigl[\Rc(\yh)\bigr],
\qquad \Rc(\yh)=\tfrac12\bigl\|\zt(\yh)-\yh\bigr\|_2^2,
\label{eq:casper}
\end{align}
where $\zt$ is the $\casper_k$ surrogate projection ($\gamma=0$). As a control we also consider
restoration towards the masked distribution, $\Rm(\yh)=\tfrac12\|\zm-\yh\|_2^2$.

\paragraph{The $\casper_k$ projection (recalled).} Let $|S|\ge k$, $T=T^\star$,
$m_T=\min_{t\in T}\yh_t$, $M_T=\max_{j\notin T}\yh_j$, and $e=\min(\epsilon,m_T)$. With
\[
L=\{t\in T:\yh_t\le M_T+e\},\qquad H=\{j\notin T:\yh_j\ge m_T-e\},\qquad B=L\cup H,
\]
\[
u=\frac{1}{|B|}\Bigl(\sum_{b\in B}\yh_b+|H|\,e\Bigr),\qquad
\zt_L=u,\quad \zt_H=u-e,\quad \zt_i=\yh_i\ (i\notin B),
\]
and $\zt=\yh$ if $L=\emptyset$. The surrogate guarantees
$\min_{t\in T}\zt_t\ge\max_{j\notin T}\zt_j+e$ (Theorem~\emph{Satisfaction}, Sec.~3.1), so
$\Topk(\zt)=T\subseteq S$. If $|S|<k$, $\zt$ is the face projection: $\zt_a=\yh_a+F/|S|$ for
$a\in S$ and $\zt_j=0$ for $j\in\Sb$.

\begin{definition}[$k$-compliance]
$\yh$ is \emph{$k$-compliant} if $\Topk(\yh)\subseteq S$. It is \emph{$\epsilon$-separated
$k$-compliant}, written $\yh\in\Ck$, if $|S|\ge k$ and $m_T>M_T+e$: the raw top-$k$ set is exactly
the $k$ best allowed values, separated from all others by more than $e$. Clearly $\Ck$ is contained
in the $k$-compliant set, and differs from it only by a margin of width $e\le\epsilon$.
\end{definition}

\section{Zero sets: restoration switches off, the penalty never does}

\begin{proposition}[Restoration vanishes exactly on the target set]\label{prop:zero}
Let $|S|\ge k$ and $\epsilon>0$.
\begin{enumerate}
\item[(i)] If $\yh\in\Ck$ then $\zt=\yh$ on a neighbourhood of $\yh$; hence $\Rc=0$ and
$\nabla_\ell\Rc=0$ there.
\item[(ii)] If $\yh$ is not $k$-compliant then $\Rc(\yh)>0$.
\end{enumerate}
\end{proposition}
\begin{proof}
(i) $m_T>M_T+e$ means $L=\emptyset$, and the strict inequality persists under small perturbations
of $\yh$ (the order statistics and $e$ are continuous, and $T$ is locally constant when the
allowed values are ranked strictly, which holds generically; with ties use the deterministic
tie-breaking rule). So $\zt\equiv\yh$ locally and $\Rc\equiv0$ locally.
(ii) If some $j\in\Sb$ lies in $\Topk(\yh)$, then some $t\in T$ lies outside $\Topk(\yh)$, so
$\yh_t\le\yh_j\le M_T$ and $m_T\le M_T$. The output instead satisfies
$\min_T\zt\ge\max_{\notin T}\zt+e>\max_{\notin T}\zt$, so $\zt\neq\yh$.
\end{proof}

\begin{proposition}[The penalty and masked restoration vanish nowhere]\label{prop:never}
For every $\yh\in\Delta^\circ$:
\begin{enumerate}
\item[(i)] $\displaystyle\frac{\partial F}{\partial\ell_i}=\yh_i\bigl(\1[i\in\Sb]-F\bigr)$; in
particular $\partial F/\partial\ell_j=\yh_jP_S>0$ for $j\in\Sb$, so $\nabla_\ell F\neq0$, and
$\inf F=0$ is not attained.
\item[(ii)] Writing $\yh_j=F w_j$ for $j\in\Sb$ (so $w\in\Delta(\Sb)$),
$\Rm(\yh)=\tfrac12F^2\bigl(\|w\|^2+\|\zm\|^2\bigr)>0$, and lowering all forbidden logits by a common
amount decreases $\Rm$ strictly; hence $\nabla_\ell\Rm\neq0$.
\item[(iii)] If $|S|<k$, $\Rc(\yh)=\tfrac12\bigl(\sum_{j\in\Sb}\yh_j^2+F^2/|S|\bigr)>0$:
$\casper_k$'s face projection has the same (empty) zero set in $\Delta^\circ$.
\end{enumerate}
\end{proposition}
\begin{proof}
(i) $\partial\yh_m/\partial\ell_i=\yh_m(\1[m=i]-\yh_i)$; sum over $m\in\Sb$.
(ii) $\zm_j-\yh_j=-Fw_j$ for $j\in\Sb$ and $\zm_a-\yh_a=\zm_a-(1-F)\zm_a=F\zm_a$ for $a\in S$. A
common shift $\ell_j\mapsto\ell_j-c$ ($j\in\Sb$) keeps $w$ and $\zm$ fixed and strictly decreases
$F$, hence $\Rm$; so the directional derivative of $\Rm$ along this direction is negative.
(iii) Direct substitution of the face projection.
\end{proof}

\begin{remark}
$\Rc$ targets $k$-compliance, the property that matters when the policy samples from its top-$k$
set with the shield removed. $F$ and $\Rm$ target $F=0$, which is stricter (it also suppresses
forbidden values far below the top-$k$ set, which can never be sampled) and unattainable with finite
logits.
\end{remark}

\section{Consequences for the optimisation}

\begin{corollary}[No bias at compliant optima]\label{cor:stationary}
Consider a tabular (logit) parametrisation $\theta=\{\ell(s)\}_s$ and a state distribution with
positive mass on some state with $S(s)\neq[h]$. Let $\theta^\star$ be a stationary point of
$\mathcal L_{\mathrm{task}}$ such that $\yh_{\theta^\star}(s)\in\Ck$ for every state with
$|S(s)|\ge k$ and $S(s)\neq[h]$. Then:
\begin{enumerate}
\item[(i)] if every constrained state has $|S(s)|\ge k$, $\theta^\star$ is a stationary point of
$\mathcal J_{\casper}$ for every $\lambda_r\ge0$;
\item[(ii)] $\theta^\star$ is a stationary point of $\mathcal J_{\mathrm{pen}}$ for no $\lambda>0$,
and of $\lambda_r\,\mathbb E_s[\Rm]$-regularised training for no $\lambda_r>0$.
\end{enumerate}
\end{corollary}
\begin{proof}
(i) By Proposition~\ref{prop:zero}(i), $\nabla_{\ell(s)}\Rc=0$ at every constrained state, and
$\Rc\equiv0$ at unconstrained ones. (ii) In the tabular case $\nabla_{\ell(s)}\mathcal J_{\mathrm{pen}}
=\nabla_{\ell(s)}\mathcal L_{\mathrm{task}}+\lambda\,\rho(s)\nabla_\ell F=\lambda\,\rho(s)\nabla_\ell F
\neq0$ by Proposition~\ref{prop:never}(i), where $\rho(s)>0$ is the state weight; likewise for $\Rm$
by Proposition~\ref{prop:never}(ii).
\end{proof}

Since masking makes $\mathcal L_{\mathrm{task}}$ invariant to the forbidden logits
(Lemma~\ref{lem:mask}), compliant task-optima exist whenever masked task-optima do: lower the
forbidden logits below the $k$-th allowed one. Corollary~\ref{cor:stationary} says $\casper_k$'s
restoration keeps those optima, at any weight, while the penalty moves every solution by an amount
that grows with $\lambda$. This is the source of the penalty's success--compliance trade-off curve.

\begin{proposition}[The penalty sharpens the executed policy]\label{prop:sharpen}
Let $\ell'=\ell-\eta\lambda\nabla_\ell F$. For $a\in S$, to first order in $\eta$,
\[
\log\zm_a(\ell')-\log\zm_a(\ell)=\eta\lambda\,F(1-F)\Bigl(\zm_a-\sum_{b\in S}(\zm_b)^2\Bigr),
\]
and the entropy of the masked policy changes by
\[
\Delta\Ent(\zm)=-\eta\lambda\,F(1-F)\operatorname{Cov}_{a\sim\zm}\bigl(\log\zm_a,\ \zm_a\bigr)\ \le\ 0,
\]
with equality iff $\zm$ is uniform on $S$. The same holds for the executed top-$k$ policy
$q^{\mathrm{mask}}$, with $S$ replaced by $T^\star$ and $F(1-F)$ by $F\sum_{b\in T^\star}\yh_b$.
\end{proposition}
\begin{proof}
By Proposition~\ref{prop:never}(i), $\Delta\ell_a=\eta\lambda F\yh_a=\eta\lambda F(1-F)\zm_a$ for
$a\in S$. Since $\log\zm_a=\ell_a-\log\sum_{b\in S}e^{\ell_b}$,
$\Delta\log\zm_a=\Delta\ell_a-\sum_b\zm_b\Delta\ell_b$, which gives the first claim. For the entropy,
$\Delta\Ent=-\sum_a\zm_a\log\zm_a\,\Delta\log\zm_a$ (the term $\sum_a\zm_a\Delta\log\zm_a$ vanishes),
which equals the stated covariance. $\log z$ and $z$ are comonotone, so the covariance is
non-negative by Chebyshev's sum inequality, and zero iff $\zm$ is constant on $S$.
\end{proof}

The penalty therefore keeps pushing the executed policy towards determinism at every state where a
forbidden value has mass, and $F>0$ always (Proposition~\ref{prop:never}). This is a
premature-commitment pressure, consistent with the extra locked-in seeds and lower task success of
the penalty in the experiments.

\begin{proposition}[Restoration is local and mass-preserving]\label{prop:local}
Let $|S|\ge k$ and assume $e=\epsilon$ (i.e.\ $\epsilon\le m_T$, so $e$ is locally constant). Then
\[
\nabla_{\yh}\Rc=\yh-\zt\quad\text{on }B,\qquad \nabla_{\yh}\Rc=0\quad\text{off }B,\qquad
\sum_{b\in B}(\yh_b-\zt_b)=0 .
\]
A descent step moves $\yh_L$ up towards $u$ and $\yh_H$ down towards $u-\epsilon$, and changes
nothing else.
\end{proposition}
\begin{proof}
Let $J=\partial\zt/\partial\yh$. For $i\notin B$, $\zt_i=\yh_i$ and $u$ does not depend on $\yh_i$, so
$J_{mi}=\1[m=i]$ and $((I-J)^\top(\yh-\zt))_i=0$. On $B$, $J_{B,B}=\frac1{|B|}\mathbf 1\mathbf
1^\top$, so $((I-J)^\top(\yh-\zt))_B=(\yh-\zt)_B-\overline{(\yh-\zt)_B}\,\mathbf 1$. Finally
$\sum_B\zt=|L|u+|H|(u-\epsilon)=|B|u-|H|\epsilon=\sum_B\yh$, so the mean vanishes, and
$\nabla_{\yh}\Rc=(I-J)^\top(\yh-\zt)=\yh-\zt$ on $B$.
\end{proof}

Comparing Propositions~\ref{prop:never}(i) and~\ref{prop:local}: the penalty's gradient touches every
value in proportion to its probability and never switches off; restoration's touches only the
forbidden values that compete for a top-$k$ slot ($H$) and the weakest top-$k$ allowed values ($L$),
moves mass only among them, and vanishes once $\yh\in\Ck$. While it is active, it raises the weakest
members of $T$, which flattens rather than sharpens the executed policy inside $T$.

\section{Predictions and empirical check}

The analysis predicts not only that $\casper_k$'s restoration should beat the penalty, but also when
the difference should disappear: whenever the target has no switch-off region.
Table~\ref{tab:check} compares the predictions with the NAVIX experiments (PPO with top-$k$ sampling;
every regulariser's weight tuned over three values with the same pre-declared rule; fresh-seed
confirmation, 32 seeds).

\begin{table}[h]
\centering
\small
\begin{tabular}{@{}>{\raggedright\arraybackslash}p{4.4cm}>{\raggedright\arraybackslash}p{3.6cm}>{\raggedright\arraybackslash}p{6.3cm}@{}}
\toprule
Regulariser / regime & Zero set in $\Delta^\circ$ & Observed (task success / violation)\\
\midrule
$\Rc$, $|S|\ge k$ ($k=2$) & $\Ck$ (switches off) &
Dyn.-Obstacles: $1.00/0.08$ and $0.81/0.06$; best trade-off curve in every setting\\
$F$ (penalty), $k=2$ & empty & Dyn.-Obstacles: $0.44/0.22$ (both settings); success falls
monotonically with $\lambda$\\
$\Rm$ (masked target), $k=2$ & empty & frozen agents (success $0$ on Dyn.-Obstacles, $\le0.05$ on
LavaGap) at the selected weight; $\approx0.3$ less success than $\Rc$ at matched violation\\
$\Rc$ with $|S|<k$ ($k=3$, face) & empty (Prop.~\ref{prop:never}(iii)) & identical to $\Rm$:
violation $0.018/0.021/0.022$ for $\Rc/\Rm/$mask${}+\Rc$, success $1.00$ for all\\
\bottomrule
\end{tabular}
\caption{The zero-set property separates the regularisers exactly where the experiments do.}
\label{tab:check}
\end{table}

All formulas in Propositions~\ref{prop:zero}--\ref{prop:local} were checked numerically against
the implementation (\texttt{ncasper/\allowbreak shields.py}), on 2{,}000--3{,}000 random states
each in float64:
the gradient of $F$ and the first-order change of $\log\zm$ match to $\le3\times10^{-16}$;
$\nabla\Rc$ is exactly $0$ on $\Ck$ and nonzero on every non-compliant state; and $\nabla_{\yh}\Rc$
equals $\yh-\zt$ to $1.1\times10^{-16}$ and vanishes off $B$.

\section{Scope}

The results are statements about the geometry of the regularisers (zero sets, stationary points,
first-order effects on the executed policy). They explain why the penalty trades task success for
compliance and why $\casper_k$'s restoration need not, and they predict where the two coincide.
They do not by themselves guarantee better final performance in reinforcement learning, which also
depends on exploration, non-convexity and parameter sharing (Corollary~\ref{cor:stationary} is stated
for the tabular parametrisation; with a shared network, the penalty's gradient is generically not
annihilated by the parameter Jacobian, but this is not guaranteed). Proposition~\ref{prop:local}
assumes $e=\epsilon$; when $m_T<\epsilon$ the margin depends on $m_T$ and an extra term appears on
$\argmin_T\yh$.

\end{document}
